% ============================================================
%  Part I — Pronouns
% ============================================================
\gpart{I}{Pronouns}

% ------------------------------------------------------------
\gsec{Personal pronouns}

\begin{casetbl}{}
  Number & Person & English & Nominativ & Akkusativ & Dativ & Genitiv \\
  \SetCell[r=6]{c,m} Singular & 1st & I            & ich & mich & mir   & meiner \\
                              & 2nd & you \eng{inf.} & du  & dich & dir   & deiner \\
                              & 2nd & you \eng{formal} & Sie & Sie & Ihnen & Ihrer \\
                              & 3rd & he           & er  & ihn  & ihm   & seiner \\
                              & 3rd & she          & sie & sie  & ihr   & ihrer \\
                              & 3rd & it           & es  & es   & ihm   & seiner \\
  \SetCell[r=3]{c,m} Plural   & 1st & we           & wir & uns  & uns   & unser \\
                              & 2nd & you \eng{inf. pl.} & ihr & euch & euch & euer \\
                              & 3rd & they         & sie & sie  & ihnen & ihrer \\
\end{casetbl}

\begin{gnote}
\textbf{Sie} (formal you) is the same word for one person or many, and it always
takes plural verb endings: \textit{Sie sind}. Capitalised mid-sentence it is
\emph{you}; lower-case \textit{sie} is \emph{she} or \emph{they}.
\end{gnote}

\begin{gnote}
The genitive column is here for completeness only. Genitive personal pronouns
are archaic and survive in fixed phrases (\textit{Wir gedenken ihrer}). In
speech use \textit{von + Dativ}. The alternative forms \textit{unserer} and
\textit{eurer} also occur.
\end{gnote}

\tcap{Watch out for \textit{ihr}}

\begin{flattbl}{colspec={X[0.7,l] X[1,l] X[2.3,l]}}
  ihr & nominative & \emph{you} (plural, informal) — \textit{ihr seid} \\
  ihr & dative     & \emph{to her} — \textit{ich gebe ihr das Buch} \\
  ihr & possessive & \emph{her} / \emph{their} — \textit{ihr Hund} \\
  Ihr & possessive & \emph{your} (formal) — \textit{Ihr Hund} \\
\end{flattbl}

% ------------------------------------------------------------
\gsec{Reflexive pronouns}

Used when the subject and the object are the same person:
\textit{ich wasche \textbf{mich}} — \emph{I wash myself}.

\begin{gtbl}{colspec={X[0.9,l] X[1.4,l] X[1,c] X[1,c]}}
  Person & Subject & Akkusativ & Dativ \\
  1st sg & ich & mich & mir  \\
  2nd sg & du  & dich & dir  \\
  3rd sg & er / sie / es & sich & sich \\
  formal & Sie & sich & sich \\
  1st pl & wir & uns  & uns  \\
  2nd pl & ihr & euch & euch \\
  3rd pl & sie & sich & sich \\
\end{gtbl}

\begin{gnote}
Only the 1st and 2nd person singular distinguish accusative from dative
(\textit{mich}/\textit{mir}, \textit{dich}/\textit{dir}). Everything else is
identical in both cases.
\end{gnote}

\tcap{Accusative or dative?}

\begin{flattbl}{colspec={X[1,l] X[2.3,l]}}
  Akkusativ & no other object: \textit{Ich wasche \textbf{mich}.} \\
  Dativ     & a direct object is already there: \textit{Ich wasche \textbf{mir} die Hände.} \\
\end{flattbl}

\tcap{Common reflexive verbs}

\begin{flattbl}{colspec={X[1.15,l] X[1,l] X[1.15,l] X[1,l]}}
  sich freuen        & to be glad       & sich setzen      & to sit down \\
  sich erinnern      & to remember      & sich beeilen     & to hurry \\
  sich interessieren & to be interested & sich fühlen      & to feel \\
  sich vorstellen    & to imagine \eng{dat.} & sich anziehen & to get dressed \\
  sich befinden      & to be located    & sich entscheiden & to decide \\
\end{flattbl}

% ------------------------------------------------------------
\gsec{Relative pronouns}

Identical to the definite article except for the five highlighted forms.

\begin{gendertbl}{}
  Case & maskulin & feminin & neutrum & Plural \\
  Nominativ & der & die & das & die \\
  Akkusativ & den & die & das & die \\
  Dativ     & dem & der & dem & \en{denen} \\
  Genitiv   & \en{dessen} & \en{deren} & \en{dessen} & \en{deren} \\
\end{gendertbl}

\begin{gnote}
The \textbf{gender} comes from the noun you are pointing back to; the
\textbf{case} comes from the job the pronoun does inside the relative clause.
Those two decisions are completely independent of each other.
\end{gnote}

\tcap{Worked examples}

\begin{extbl}{colspec={X[1.15,l] X[1.85,l]}}
  Der Mann, \textbf{der} dort steht    & subject $\rightarrow$ nominative, masc. \\
  Der Mann, \textbf{den} ich kenne     & direct object $\rightarrow$ accusative, masc. \\
  Der Mann, \textbf{dem} ich helfe     & \textit{helfen} takes the dative \\
  Die Frau, \textbf{deren} Auto neu ist & possession $\rightarrow$ genitive, fem. \\
  Die Leute, mit \textbf{denen} ich rede & \textit{mit} $\rightarrow$ dative plural \\
\end{extbl}

\begin{gnote}
The verb goes to the \textbf{end} of a relative clause, and the clause is always
fenced off with commas. Use \textit{was} instead of a der-form after
\textit{alles, nichts, etwas, vieles} and after a whole preceding clause.
\end{gnote}
